\documentclass[11pt]{article}
\usepackage{amsmath,amssymb,hyperref}
\title{Who can audit an old duty after a change of authority?}
\author{}\date{}
\begin{document}\maketitle

\paragraph{Scope and actors.}
This is a finite strict-incentive refinement of
\href{https://github.com/eternal-november/eos-research-programmes/pull/663}
{the attributable-refusal result}. It concerns already enrolled
civil service operators with fixed tolls and deposited bonds,
not selecting the finance contract or establishing participation.
At a named public missed-duty observation \(M\), an eligible operator
\(i\) can save \(d_i\ge0\) by refusing rather than performing, while
an unable operator is compliant. There is a possible unilateral
deviation history consistent with \(M\) attributable to each checked
subject \(i\). For example, every other missed duty could be genuinely
unavailable if at least \(|M|-1\) such failures are admitted. That
condition is not assumed for cases outside the specified observation
class. The intended reference holds the operator's information and
pre-existing financial terms fixed. It is not a simultaneous
coalitional equilibrium or a claim about real financial enforcement.

The settlement service commits in advance to a lottery over lawful
audit allocations. An admissible unit audit names the earlier duty,
its eligibility and historical holder rather than just the new holder
of its equipment. It correctly distinguishes refusal from inability,
can settle before the financial deadline, and charges at most the
existing bond \(B_i>0\) only after attributable refusal is proved.
No audit can confer physical access or authorize a late rescue.
The operator does not learn the random audit draw before deciding,
cannot influence the verifier, and gains \(d_i\) if unaudited. A
qualified audit is certain to detect that modeled refusal; lower
disclosure or reliability needs a new edge/evidence contract. Under
these assumptions, strict compliance at every such eligible
unilateral deviation is equivalent to \(p_i>r_i=d_i/B_i\), where
\(p_i\) is conditional on the public observation and on the appropriate
authority state. The law of no punishment for honest inability
holds separately in every role. A profitable deviation not in the
specified family may invalidate the incentive result.

\paragraph{Known late authority: exact feasibility and a realizable lottery.}
Represent one audit as an edge from a subject to a verifier with
integer capacity \(k_v\); edges require both permission to use the
historical evidence and time to produce it. Each draw is an actual
matching with at most one audit of each subject. Let \(N(S)\) denote
the verifier neighbourhood of subjects \(S\). The exact largest
common additive slack over the incentive thresholds is
\[
 \epsilon_* =
 \min\!\left\{ \min_i (1-r_i),\quad
    \min_{\varnothing\ne S\subseteq M}
   \frac{\sum_{v\in N(S)} k_v - \sum_{i\in S}r_i}{|S|}\right\}.
\]
Positive slack is necessary and sufficient for the strict contract.
Necessity follows from the unit individual bounds and checking the
expected total audits across each cut. For sufficiency, use
\(r_i+\epsilon\) with \(0<\epsilon\le\epsilon_*\). A flow from the
subjects of those demands through the allowed verifiers fits the
capacities by max flow/min cut. Split each verifier into \(k_v\)
unit slots. Clearing denominators gives a bipartite multigraph with
maximum degree \(Q\), the common denominator. Pad both sides and
fill deficits with dummy edges to make it \(Q\)-regular. Repeatedly
remove a perfect matching, discarding dummy edges in the deployed
schedule. A uniform draw among the resulting \(Q\) schedules gives
the desired rational marginals, and no schedule overspends a verifier.
Max flow yields the same upper slack for small or negative
\(\epsilon_*\): at a negative proposed slack the audit lower bounds
are \(\max(0,r_i+\epsilon)\). In every cut, retaining only the positively
constrained subjects can only shrink the required neighbourhood.
The exhaustive cut and decomposition code is a reference
implementation; its length can depend on the denominators.

For \(r=(3/5,3/5,2/5)\), two single-slot verifiers have more total
capacity than total threshold. Initially the first two subjects
share just one admissible verifier, so strict compliance is impossible
and the common slack is \(-1/10\). A successor arrangement granting
the second verifier valid historical-audit power for the middle duty
changes no capacity and moves the slack to \(2/15\). The explicit
lottery delivers probabilities \(2/3,2/3,7/15\), all strict.
A nominal appointment after the settlement cut changes no edge. The
native entitlement example in the prior study motivates the named-old
duty premise; it does not establish any proposed verifier or lottery.

\paragraph{A different cut: choose the subjects before recognizing the role.}
Suppose the audit subjects must be chosen irreversibly while several
verifier graphs are compatible with the public record. Only after that
choice is the actual graph known to the assignment layer, in time to
make a lawful assignment. A valid draw must be matchable in
\emph{every} possible graph, but can use a different particular matching
in each. The exact marginal region is the convex hull of indicators
of those common matchable sets. Intersecting the edge lists is too
conservative: a single subject eligible for verifier A in one role
and only B in another is always auditable if allocation waits.
Intersecting separately achievable marginal vectors can instead
be too permissive.

With only two graphs, separately achievable identical marginals do
suffice for this timing pattern. Route flow from the first role's
verifier slots through each subject's unit-capacity edge, to the
second role's slots. The two given fractional assignments fit together
at the shared subject edge; an integer-capacity circulation
(add a return edge from sink to source) decomposes into integral
circulations. Each integral draw supplies the same subjects and a
valid assignment to each graph. This is a direct flow instance of the
classical two-transversal-matroid intersection principle, not a new
general combinatorial theorem.

Three graphs can already fail. In each graph exactly one different
pair shares one verifier, with a separate second verifier for the
third subject. Each recognized role supports inclusion \(1/2\) for
each of the three subjects, and the first two graphs also have a
common lottery with those marginals. Across all three, no pair of
subjects is matchable in every role. A draw therefore audits at most
one, so three thresholds \(49/100\) cannot all be strictly exceeded.
For a family with one such conflicting-pair role for every edge of an
arbitrary finite graph \(G\) and separate slots for all other subjects,
the feasible draws are exactly the independent sets of \(G\).
The maximum common inclusion probability equals \(1/\chi_f(G)\):
normalize independent-set covering weights to a lottery for one
direction, and scale a lottery with minimum marginal \(p>0\) by
\(1/p\) to get covering weights for the other. In a five-cycle each
single role permits one half, while early choice permits only two
fifths. A contract can restore the higher threshold by delaying
subject selection until the applicable role is known, provided it
can still obtain evidence before the cutoff. This intervention is
different from learning who holds a nominal title, and can fail when
an immutable evidence window demands earlier selection.

\paragraph{Comparators, checks and what does not transfer.}
The strong direct comparator enumerates all legal deterministic
assignments and solves for a lottery, not an aggregate-capacity-only
policy. It agrees on 576 known-graph cases, including 183 with a
constructed strict lottery. Common subject versus edge sets, two and
three roles, all 64 four-vertex graphs and the five-cycle were checked
against explicit feasible draws and linear programs. Selected strict
capacity obstructions were checked in Lean. The broader
\href{https://economics.sas.upenn.edu/sites/default/files/filevault/SSRN%2016_009%20third%20version.pdf}
{costly-verification literature} and
\href{https://www.sciencedirect.com/science/article/pii/S0167506008708173}
{matroid intersection} supply comparison structure. The contribution
here is the governance timing and historical-authority interpretation
with a concrete discriminator, not novelty of Hall, bipartite matching
or fractional coloring. An auditor whose legal or physical scope
is not established, a manipulable draw, a binding participation
constraint or shared penalty exposure needs another analysis. No
architecture or independent R5 acceptance is claimed.

\end{document}
